%% kapa: exam=1 round=34 year=2023 period=1 form=A subject=economics subject_ko=경제학원론 seq=33 number=73
%% source: Q-Net 감정평가사 기출문제 (cst00309) artlSeq=5211789 file=raw/1008/5211789/2023년도 제34회 감정평가사 1차시험 1교시.pdf
%% conversion: kordoc 4.24.51 -> tools/gen1.py (auto); review=proofread-sonnet; flags=none
\begin{question}{73}
인플레이션에 관한 설명으로 옳지 않은 것은?
\begin{choices}
  \item 메뉴비용은 인플레이션에 맞춰 가격을 변경하는 데에 발생하는 각종 비용을 말한다.
  \item 디스인플레이션(disinflation) 상황에서는 물가상승률이 감소하고 있지만 여전히 물가는 상승한다.
  \item 초인플레이션은 극단적이고 장기적인 인플레이션으로 통제가 어려운 상황을 말한다.
  \item 구두창비용은 인플레이션에 따라 발생하는 현금 관리 비용을 말한다.
  \item 디플레이션은 인플레이션이 진행되는 상황에서 경제가 침체하는 상황을 말한다.
\end{choices}
\end{question}
\answer{5}
